% ANEXO A - Documentação de APIs e Referências Técnicas
\chapter{Documentação de Referências Técnicas}
\label{anexo:a}

Este anexo apresenta detalhes técnicos exaustivos sobre as interfaces de integração, padrões de resposta e estruturas de dados do ecossistema Oficina Master.

\section{Dicionário de Dados do Hub Master}

Abaixo detalham-se as propriedades das entidades centrais de orquestração.

\subsection{Entidade: Tenant (Configurações de Instância)}
\begin{longtable}{lll}
    \toprule
    \textbf{Campo} & \textbf{Tipo} & \textbf{Descrição} \\
    \midrule
    \texttt{uuid} & char(36) & Identificador único global (Primary Key). \\
    \texttt{name} & varchar(255) & Nome fantasia do tenant no sistema. \\
    \texttt{database} & varchar(60) & Nome lógico do esquema no servidor MySQL. \\
    \texttt{url\_banco} & varchar(255) & Host/IP do servidor de banco de dados. \\
    \texttt{username} & varchar(60) & Usuário com permissões de owner no banco. \\
    \texttt{password} & text & Senha criptografada (AES-256). \\
    \texttt{status} & enum & Estado operacional (ativo, bloqueado, desativado). \\
    \texttt{licenca\_valida\_ate} & date & Cache da data de expiração da licença atual. \\
    \bottomrule
\end{longtable}

\section{Especificação Completa de Endpoints (API Master)}

A comunicação entre o \textit{frontend} e o \textit{backend} é realizada via HTTPS utilizando o padrão JSON.

\subsection{Autenticação e Segurança}
\begin{itemize}
    \item \texttt{POST /api/auth/login}: Autenticação de usuários Master (Retorna Bearer Token);
    \item \texttt{GET /api/auth/me}: Retorna dados do perfil logado e permissões;
    \item \texttt{POST /api/auth/logout}: Invalida o token de acesso atual.
\end{itemize}

\subsection{Gestão de Tenants e Infraestrutura}
\begin{itemize}
    \item \texttt{GET /api/empresas}: Listagem paginada com filtros por status e nome;
    \item \texttt{POST /api/empresas/migrate/\{id\}}: Dispara o \textit{job} de inicialização forçada de banco;
    \item \texttt{GET /api/database/tables/\{uuid\}}: Lista tabelas e contagem de registros de um banco isolado;
    \item \texttt{POST /api/database/test-connection}: Valida credenciais antes da persistência.
\end{itemize}

\section{Padrão de Resposta JSON (Standard API Response)}

Todas as respostas da API seguem o esquema estruturado abaixo para facilitar o consumo pelo \textit{frontend}:

\begin{lstlisting}[language=json, caption={Exemplo de Resposta de Sucesso}]
{
  "status": true,
  "code": "SUCCESS",
  "data": { ... },
  "message": "Operação realizada com sucesso",
  "meta": {
    "current_page": 1,
    "total": 150
  }
}
\end{lstlisting}

\section{Consultas de Telemetria SQL (Referência)}

Para a monitoria dos bancos de dados, o sistema utiliza consultas de inspeção de esquema:

\begin{lstlisting}[language=sql, caption={Query de Inspeção de Tamanho de Tabelas}]
SELECT 
    table_name AS "Tabela",
    table_rows AS "Registros",
    round(((data_length + index_length) / 1024 / 1024), 2) AS "Tamanho (MB)"
FROM information_schema.TABLES
WHERE table_schema = 'db_cliente_uuid'
ORDER BY (data_length + index_length) DESC;
\end{lstlisting}
